\chapter{The formalization target}

\section{Scope and source}

This project formalizes the combinatorial proof of the density Hales--Jewett theorem due to
Dodos, Kanellopoulos, and Tyros \cite{DKT}.  Their proof uses only the uniform probability
measure on finite word spaces.  The purpose of this blueprint is twofold: it records the complete
mathematical dependency graph of that proof, and it refines arguments which are compressed or
omitted in the paper into units suitable for Lean.

The most substantial refinement is Chapter~\ref{chap:gr}.  Proposition~2 of \cite{DKT} is quoted
as a consequence of the Graham--Rothschild theorem.  Here it is proved from the coloring
Hales--Jewett theorem by a finite-unions argument and a block-canonization argument.  We also spell
out the final density-increment iteration, the estimates hidden in several averaging steps, and the
deduction of arithmetic progressions.

The paper writes $[k]=\{1,\ldots,k\}$ and $[k]^n$ for words.  Lean will instead use a finite
alphabet $\alpha$ and a finite coordinate type $\iota$, with word space $\iota\to\alpha$.  In the
quantitative parts we specialize to $\alpha=\operatorname{Fin} k$ and
$\iota=\operatorname{Fin} n$.  This agrees with the existing mathlib API and avoids continual
transport between vectors and functions.

\begin{remark}[Status convention]
A declaration annotated by \texttt{\textbackslash mathlibok} already exists in mathlib.  A project
declaration without \texttt{\textbackslash leanok} is part of the proposed implementation.  At the
time of writing, only the two final challenge statements exist in the project, and both have
placeholder proofs.
\end{remark}

\section{Main statements}

\begin{theorem}[Density Hales--Jewett]
\label{thm:dhj}
\lean{Combinatorics.Line.exists_of_density}
Let $\alpha$ be a finite alphabet and let $\delta>0$.  There is a natural number
$N=N(\card\alpha,\delta)$ such that, whenever $n\geq N$ and
$A\subseteq(\operatorname{Fin}n\to\alpha)$ satisfies
\[
  \delta\,\card\alpha^n\leq\card A,
\]
there is a combinatorial line $\ell$ such that $\ell(a)\in A$ for every $a\in\alpha$.
\end{theorem}

The cases $\card\alpha\leq 1$ are elementary and will be separated before invoking the induction
on the alphabet size.  This lets the main argument retain the paper's standing assumption
$k\geq2$.

\begin{theorem}[Szemer\'{e}di's theorem for finite intervals]
\label{thm:szemeredi}
\lean{Combinatorics.ArithmeticProgression.exists_of_density_nat}
\uses{thm:dhj}
For every $k\in\N$ and every $\delta>0$ there is a natural number $N=N(k,\delta)$ such
that, whenever $n\geq N$ and $A\subseteq\{0,\ldots,n-1\}$ satisfies
$\delta n\leq\card A$, the set $A$ contains all the terms of a nonconstant arithmetic
progression of length $k$.
\end{theorem}

The proof of Theorem~\ref{thm:dhj} occupies Chapters~\ref{chap:geometry}--\ref{chap:iteration}.
The short digital transfer proving Theorem~\ref{thm:szemeredi} is given in
Chapter~\ref{chap:ap}.


\chapter{Finite word geometry}
\label{chap:geometry}

\section{Lines and subspaces}

\begin{definition}[Words]
\label{def:word}
For finite types $\iota$ and $\alpha$, a word over $\alpha$ with coordinate set $\iota$ is a
function $w:\iota\to\alpha$.  The full word space is denoted
$\mathbb W(\alpha,\iota)=\iota\to\alpha$.
\end{definition}

\begin{definition}[Combinatorial subspace]
\label{def:subspace}
\lean{Combinatorics.Subspace}
\mathlibok
Let $\eta$ be a finite direction type.  A combinatorial subspace
$V:\operatorname{Subspace}\eta\,\alpha\,\iota$ is specified by a map
\[
  V.\operatorname{idxFun}:\iota\longrightarrow\alpha\sqcup\eta
\]
such that every $e\in\eta$ occurs in the right summand.  It evaluates to an embedding
\[
  V:(\eta\to\alpha)\longrightarrow(\iota\to\alpha)
\]
by using a fixed letter at coordinates in the left summand and the input letter $x(e)$ at
coordinates labelled by $e$.  Its range is the corresponding $\card\eta$-dimensional word
subspace.
\end{definition}

\begin{definition}[Combinatorial line]
\label{def:line}
\lean{Combinatorics.Line}
\mathlibok
A line $\ell:\operatorname{Line}\alpha\,\iota$ is specified by a map
$\iota\to\operatorname{Option}\alpha$ which is \texttt{none} at least once.  The coordinates
labelled \texttt{none} form the nonempty variable set; the other coordinates contain fixed
letters.  Evaluation gives $\ell:\alpha\to(\iota\to\alpha)$.
\end{definition}

\begin{lemma}[Subspaces are embeddings]
\label{lem:subspace-injective}
\lean{DensityHalesJewett.Subspace.injective}
\uses{def:subspace}
If the alphabet is nontrivial, evaluation by a subspace is injective.  Consequently a subspace
with direction type $\eta$ has exactly $\card\alpha^{\card\eta}$ points.
\end{lemma}

The existing lemma states injectivity of the map which sends a subspace structure to its
evaluation function.  The implementation also needs the elementary injectivity of the evaluation
map of a fixed subspace; it follows by choosing, for each direction, a coordinate supplied by the
\texttt{proper} field.

\begin{definition}[Range and containment]
\label{def:subspace-range}
\lean{DensityHalesJewett.Subspace.range}
\uses{def:subspace}
The range of $V$ is the finite set
\[
  [V]=\{V(x):x\in(\eta\to\alpha)\}.
\]
A line or subspace is contained in $A$ if all of its evaluations belong to $A$.
For a subspace $V$, let $\Lines(V)$ denote the finite set of line structures in the ambient word
space whose ranges are contained in $[V]$.
\end{definition}

In Lean it is preferable to work with parameterized line structures rather than sets of points.
For nontrivial alphabets a line is determined by its evaluation function, so no multiplicity is
introduced.  Composition of a line in the parameter cube with $V$ provides an equivalence between
lines of $\eta\to\alpha$ and lines contained in $V$.

\begin{lemma}[Lines in a subspace]
\label{lem:lines-equiv}
\lean{DensityHalesJewett.Subspace.linesEquiv}
\uses{def:line,def:subspace-range,lem:subspace-injective}
For a nontrivial finite alphabet, composition with an $m$-dimensional subspace $V$ is an
equivalence
\[
  \operatorname{Line}\alpha\,(\operatorname{Fin}m)
  \simeq \Lines(V).
\]
In particular,
\[
  \card{\Lines(V)}=(k+1)^m-k^m
  \qquad\text{when }\card\alpha=k.
\]
\end{lemma}

\begin{proof}
A line structure is a map from the $m$ directions to $\operatorname{Option}\alpha$, with at least
one \texttt{none}.  There are $(k+1)^m$ maps in total and exactly $k^m$ constant-only maps.  For
composition, replace each occurrence of a direction $e$ in the generating word for $V$ by the
entry of the parameter line at $e$.  Injectivity and surjectivity follow by inspecting one proper
coordinate for each direction.
\end{proof}

\section{Density, restriction, and fibers}

\begin{definition}[Uniform density]
\label{def:density}
\lean{DensityHalesJewett.density}
For a finite set $A\subseteq\mathbb W(\alpha,\iota)$ define
\[
  \dens(A)=\frac{\card A}{\card\alpha^{\card\iota}}\in\R.
\]
For a subspace $V$, define the relative density
\[
  \dens_V(A)=\frac{\card{A\cap[V]}}{\card{[V]}}.
\]
\end{definition}

We will use normalized real-valued density internally.  The cardinal inequalities in the public
API are converted to density inequalities once, using the cardinality of the full function type.

\begin{definition}[Restriction of the alphabet]
\label{def:alphabet-restriction}
\lean{DensityHalesJewett.Subspace.restrictAlphabet}
\uses{def:subspace}
Given an embedding $e:\beta\hookrightarrow\alpha$ and a subspace $V$ over $\alpha$, define
\[
  V\restr\beta=\{V(e\circ x):x:\eta\to\beta\}.
\]
For the induction step $\alpha=\operatorname{Fin}(k+1)$ and
$\beta=\operatorname{Fin}k$, using the standard inclusion.
\end{definition}

This is a finite set, not generally the range of a subspace over the smaller alphabet without a
transport of the fixed letters.  All subspaces to which the operation is applied in the main proof
have fixed letters in $\operatorname{Fin}k$, so the implementation will package that side condition
and produce an actual restricted subspace when required.

\begin{definition}[Concatenation and fibers]
\label{def:fiber}
\lean{DensityHalesJewett.fiber}
Let $\iota$ and $\kappa$ be disjoint coordinate types.  Under the canonical equivalence
$(\iota\sqcup\kappa\to\alpha)\simeq(\iota\to\alpha)\times(\kappa\to\alpha)$, write
$x\concat y$ for concatenation.  For $A\subseteq\mathbb W(\alpha,\iota\sqcup\kappa)$ and
$x:\iota\to\alpha$, define
\[
  A_x=\{y:\kappa\to\alpha:x\concat y\in A\}.
\]
\end{definition}

Using sums of coordinate types is more robust than repeatedly proving identities involving
$n-l$.  Statements over $\operatorname{Fin}n$ are recovered using the standard equivalence
$\operatorname{Fin}l\sqcup\operatorname{Fin}(n-l)\simeq\operatorname{Fin}n$.

\begin{lemma}[Fiber averaging]
\label{lem:fiber-average}
\lean{DensityHalesJewett.average_density_fiber}
\uses{def:density,def:fiber}
For every finite $A\subseteq\mathbb W(\alpha,\iota\sqcup\kappa)$,
\[
  \E_{x\in\mathbb W(\alpha,\iota)}\dens(A_x)=\dens(A).
\]
The analogous identity holds when the prefix space is replaced by the range of a subspace.
\end{lemma}

\begin{proof}
Double-count the pairs $(x,y)$ for which $x\concat y\in A$, then divide by
$\card\alpha^{\card\iota+\card\kappa}$.  The subspace version uses
Lemma~\ref{lem:subspace-injective} to identify its range with a full parameter cube.
\end{proof}

\begin{lemma}[Threshold averaging]
\label{lem:threshold-average}
\lean{DensityHalesJewett.density_ge_threshold}
Let $f:X\to[0,1]$ on a nonempty finite type.  If $\E_x f(x)\geq a$ and
$0\leq b<a$, then
\[
 \dens\{x:f(x)\geq b\}\geq\frac{a-b}{1-b}.
\]
We will freely use weaker consequences such as a lower bound of $a/2$ after taking $b=a/2$.
\end{lemma}

\begin{proof}
If $H=\{x:f(x)\geq b\}$, then
$\E f\leq\dens(H)\cdot1+(1-\dens(H))b$.  Rearrangement gives the result.
\end{proof}


\chapter{The coloring input: Graham--Rothschild for lines}
\label{chap:gr}

The density argument needs a two-color Ramsey theorem for the lines of a word cube.  We prove the
exact special case required by \cite{DKT}, rather than formalizing the full parameter-sets theorem.

\section{Finite unions from Hales--Jewett}

For finite nonempty subsets $B,C\subseteq\N$, write $B<C$ if every element of $B$ is smaller than
every element of $C$.  A block sequence is a tuple $B_1<\cdots<B_t$ of nonempty sets, and
\[
  \FU(B_1,\ldots,B_t)
  =\left\{\bigcup_{i\in I}B_i:\varnothing\neq I\subseteq\{1,\ldots,t\}\right\}.
\]

\begin{theorem}[Coloring Hales--Jewett]
\label{thm:color-hj}
\lean{Combinatorics.Line.exists_mono_in_high_dimension}
\mathlibok
For finite types $\alpha$ and $\kappa$, some finite coordinate type $\iota$ has the property that
every coloring $(\iota\to\alpha)\to\kappa$ admits a monochromatic combinatorial line.
\end{theorem}

\begin{lemma}[Finite-unions focusing step]
\label{lem:fu-focus}
\lean{DensityHalesJewett.FiniteUnions.focus}
\uses{thm:color-hj}
Fix a finite color type $C$ and $q\in\N$.  Given sufficiently many consecutive reservoir blocks
and a coloring of their nonempty unions, one can replace a nonempty group of consecutive reservoir
blocks by a single new block $B$ so that the following finite profile is constant:
\[
  \left(
  \chi\left(U\cup\bigcup_{j=1}^{h}R_j(t_j)\right)
  \right)_{U\in\mathcal U}.
\]
Here $\mathcal U$ is any previously fixed family of at most $q$ unions, and each $R_j(t_j)$ is one
of finitely many prescribed choices from the $j$th reservoir packet.  In particular, substituting
any letter along the selected Hales--Jewett line leaves every color in the profile unchanged.
\end{lemma}

\begin{proof}
Partition a long reservoir into $h$ equal packets.  A word
$t=(t_1,\ldots,t_h)$ over the finite alphabet of allowed packet choices determines all the unions
appearing in the displayed profile.  Color $t$ by the entire tuple of their colors; this is a
coloring by the finite type $C^{\mathcal U}$.  Apply Theorem~\ref{thm:color-hj}.  The union of the
packet pieces at the variable coordinates is nonempty and becomes the new block $B$; fixed
coordinates are absorbed into the fixed part of the reservoir.  Monochromaticity of the line is
exactly the asserted simultaneous invariance.  Empty unions are assigned an arbitrary dummy color
when constructing the profile and are never used in the conclusion.
\end{proof}

\begin{theorem}[Finite unions theorem]
\label{thm:finite-unions}
\lean{DensityHalesJewett.FiniteUnions.exists_mono}
\uses{lem:fu-focus}
For every finite color type $C$ and every $m\in\N$, there is $L$ such that every coloring of the
nonempty subsets of $\operatorname{Fin}L$ by $C$ admits a block sequence
$B_1<\cdots<B_m$ for which $\FU(B_1,\ldots,B_m)$ is monochromatic.
\end{theorem}

\begin{proof}
We give the finite color-focusing proof.  Induct on $m$.  There is nothing to prove for $m=0$,
and $L=1$ works for $m=1$.

For the induction step, reserve enough consecutive packets to perform the focusing step once for
every member of
$\FU(B_1,\ldots,B_m)\cup\{\varnothing\}$; this is a finite list of size $2^m$.
Apply Lemma~\ref{lem:fu-focus} successively, each time coloring by the full profile already fixed.
At the end there is a new block $B_{m+1}$, later than the existing blocks, such that
\[
  \chi(U\cup B_{m+1})=\chi(U)
  \quad\text{for every }U\in\FU(B_1,\ldots,B_m),                 \tag{\(*\)}
\]
and the focusing profile also makes $\chi(B_{m+1})$ equal to this common color.  Before this final
selection, apply the induction hypothesis to the induced profile-coloring of candidate earlier
blocks.  It makes all nonempty $U$ on the right of \((*)\) have one color.  Thus every nonempty
union of $B_1,\ldots,B_{m+1}$ has that color: unions not using $B_{m+1}$ by induction, the singleton
$B_{m+1}$ by the last profile coordinate, and the remaining unions by \((*)\).

All choices are finite.  Reading the construction backwards gives a primitive recursive bound:
at each focusing call use a Hales--Jewett number with color type a finite power of $C$, and take the
sum of the required packet lengths.  The implementation will define the bound by this recursion;
no least-number construction is needed.
\end{proof}

\section{Canonizing colors of variable words}

A one-variable word of length $n$ is another presentation of a line in $[k]^n$.  Its support is
the nonempty set of coordinates occupied by the variable.

\begin{lemma}[Block canonization]
\label{lem:block-canonization}
\lean{DensityHalesJewett.GrahamRothschild.canonize}
\uses{thm:color-hj}
Let $A$ and $C$ be finite nonempty types and let $L\geq1$.  There are block lengths
$M_1,\ldots,M_L$ such that every $C$-coloring $\chi$ of the variable words of total length
$M=M_1+\cdots+M_L$ admits variable words $w_i(v)$ in the $i$th block with this property:
the color of
\[
  w_1(y_1)\concat\cdots\concat w_L(y_L),
  \qquad y_i\in A\cup\{v\},
\]
whenever at least one $y_i=v$, depends only on the support
$\{i:y_i=v\}$ and not on the fixed letters $y_i\in A$.
\end{lemma}

\begin{proof}
Choose the block lengths recursively and select the words in reverse order.  Suppose the words in
blocks $i+1,\ldots,L$ have been selected.  A constant word $u\in A^{M_i}$ is colored by its entire
finite profile: for every prefix in the earlier blocks and every choice of either a fixed letter or
the variable in each later selected block, record the color of the resulting variable word.
Contexts with no variable receive a dummy value.  The profile type is finite, so choose $M_i$ to
be a Hales--Jewett dimension for that type.  A monochromatic line in $A^{M_i}$ supplies $w_i(v)$.
Its monochromaticity says simultaneously that the original color is independent of the letter
substituted into block $i$ in every valid context.

After processing all blocks, compare two patterns with the same variable support.  Change their
fixed letters one block at a time.  At each change some variable block remains, so the relevant
context is valid and the selected profile gives equality of colors.  This proves the claim.
\end{proof}

\begin{theorem}[Graham--Rothschild, line case]
\label{thm:graham-rothschild}
\lean{DensityHalesJewett.GrahamRothschild.lines}
\uses{thm:finite-unions,lem:block-canonization,lem:lines-equiv}
For every $k\geq2$, $r\geq1$, and $m\geq1$, there is
$\GR(k,r,m)$ such that every $r$-coloring of the combinatorial lines of $[k]^n$, with
$n\geq\GR(k,r,m)$, admits an $m$-dimensional subspace $V$ all of whose lines have the same color.
\end{theorem}

\begin{proof}
Let $L$ be supplied by Theorem~\ref{thm:finite-unions} for $r$ colors and $m$ blocks, and choose
$M=M_1+\cdots+M_L$ from Lemma~\ref{lem:block-canonization}.  It suffices to prove the theorem in
dimension $M$; extra coordinates are fixed to an arbitrary letter.

Apply block canonization to the coloring of one-variable words.  It induces a well-defined coloring
$\bar\chi$ of the nonempty subsets of $\{1,\ldots,L\}$: the value on $S$ is the color of any
concatenation in which precisely the blocks in $S$ contain the variable.  By finite unions, choose
$B_1<\cdots<B_m$ such that all nonempty unions of the $B_j$ have one $\bar\chi$-color.

Fix $a_0\in[k]$.  Define an $m$-variable word $z(v_1,\ldots,v_m)$ blockwise: in block $i$, use
$w_i(v_j)$ if $i\in B_j$, and use $w_i(a_0)$ if $i$ belongs to no $B_j$.  Each $B_j$ is nonempty,
so every $v_j$ occurs and $z$ generates an $m$-dimensional subspace $V$.

Any line in $V$ is obtained by substituting into $(v_1,\ldots,v_m)$ a one-variable pattern
$(y_1,\ldots,y_m)$ with $y_j\in[k]\cup\{v\}$ and nonempty variable index set $I$.  Its support in
the $L$ canonical blocks is exactly $\bigcup_{j\in I}B_j$.  Block canonization discards the fixed
letters, and the finite-unions choice makes the resulting color independent of nonempty $I$.
Hence every line in $V$ has the same color.
\end{proof}

\begin{corollary}[The form used in the density proof]
\label{cor:gr-two-color}
\lean{DensityHalesJewett.GrahamRothschild.lines_twoColor}
\uses{thm:graham-rothschild}
For $k\geq2$ and $m\geq1$ there is $\GR(k,m)=\GR(k,2,m)$ such that, for every family
$\cL$ of lines of $[k]^n$ with $n\geq\GR(k,m)$, some $m$-dimensional subspace $V$ satisfies
either $\Lines(V)\subseteq\cL$ or $\Lines(V)\cap\cL=\varnothing$.
\end{corollary}


\chapter{Preliminary density lemmas}
\label{chap:prelim}

For $k\geq2$, write $\mathsf{DHJ}_k$ for the assertion that the density Hales--Jewett theorem is
known for the $k$-letter alphabet at every positive density.  Bounds are chosen monotone in the
density parameter by replacing any bound $B(k,\delta)$ with the maximum of the finitely many bounds
needed on a rational grid, or, more simply, by defining the least valid bound after existence has
been proved.

\begin{proposition}[Multidimensional density Hales--Jewett]
\label{prop:mdhj}
\lean{DensityHalesJewett.Subspace.exists_of_density}
\uses{thm:dhj,lem:fiber-average,lem:threshold-average}
Assume $\mathsf{DHJ}_k$.  For every $m\geq1$ and $\delta>0$ there is
$\MDHJ(k,m,\delta)$ such that every $A\subseteq[k]^n$ of density at least $\delta$, with
$n\geq\MDHJ(k,m,\delta)$, contains an $m$-dimensional subspace.
\end{proposition}

\begin{proof}
Induct on $m$.  The case $m=1$ is $\mathsf{DHJ}_k$.  Suppose the result is known in dimension
$m$, put $M=\DHJ(k,\delta/2)$, and set
\[
 \MDHJ(k,m+1,\delta)
 =M+\MDHJ\left(k,m,\frac{\delta}{2}(k+1)^{-M}\right).
\]
For $A\subseteq[k]^n$, fiber over the first $n-M$ coordinates.  By
Lemmas~\ref{lem:fiber-average} and~\ref{lem:threshold-average}, the set $B$ of prefixes $x$ with
$\dens(A_x)\geq\delta/2$ has density at least $\delta/2$.  Each $A_x$ contains a line $\ell_x$.
There are fewer than $(k+1)^M$ possible line structures, so one line $\ell$ occurs for a set $C$
of prefixes of density at least $\frac\delta2(k+1)^{-M}$.  The induction hypothesis gives an
$m$-subspace $W\subseteq C$, and $W\concat\ell$ is an $(m+1)$-subspace contained in $A$.
\end{proof}

\begin{lemma}[Uniform fibers on a subspace]
\label{lem:uniform-fibers}
\lean{DensityHalesJewett.Subspace.exists_fibers_dense}
\uses{lem:fiber-average,lem:threshold-average}
Let $k\geq2$, $m\geq1$, and $0<\varepsilon<1$.  If
$n\geq\varepsilon^{-1}k^m m$ and $A\subseteq[k]^n$ has density greater than $\varepsilon$, then
there are $l<n$ and an $m$-dimensional subspace $V\subseteq[k]^l$ such that
\[
  \dens(A_x)\geq\dens(A)-\varepsilon
  \quad\text{for every }x\in V.
\]
\end{lemma}

\begin{proof}
Set $\rho=\varepsilon/(k^m-1)$.  Start with the first free $m$-coordinate block $V_1=[k]^m$.
If every fiber above $V_1$ has density at least $\dens(A)-\varepsilon$, stop.  Otherwise, since
their average is $\dens(A)$, one fiber has density at least $\dens(A)+\rho$: if one value is below
$\dens(A)-\varepsilon$, the remaining $k^m-1$ values must compensate by at least $\rho$.

Fix that exceptional prefix and repeat on the next $m$ coordinates.  After $s$ failures the chosen
fiber has density at least $\dens(A)+s\rho$.  Since density is at most one, there are at most
$\lfloor\rho^{-1}\rfloor$ failures.  The hypothesis on $n$ leaves another block of $m$ coordinates,
so the procedure stops and the final free block is the desired $V$.  Lean will formulate the
iteration with a fuel parameter $\lfloor\rho^{-1}\rfloor+1$.
\end{proof}

\begin{corollary}[Restricted-alphabet subspace]
\label{cor:restricted-subspace}
\lean{DensityHalesJewett.Subspace.exists_restrictAlphabet_subset}
\uses{prop:mdhj,lem:uniform-fibers,def:alphabet-restriction}
Assume $\mathsf{DHJ}_k$.  For every $m\geq1$ and $\delta>0$ there is
$\MDHJ^*(k,m,\delta)$ such that every $A\subseteq[k+1]^n$ of density at least $\delta$, with
$n\geq\MDHJ^*(k,m,\delta)$, admits an $m$-dimensional subspace $V$ satisfying
$V\restr[k]\subseteq A$.
\end{corollary}

\begin{proof}
Put $M=\MDHJ(k,m,\delta/2)$ and take
\[
  \MDHJ^*(k,m,\delta)
  =\left\lceil(\delta/2)^{-1}(k+1)^M M\right\rceil.
\]
Lemma~\ref{lem:uniform-fibers} gives an $M$-subspace $W$ in an initial coordinate block such that
every fiber above $W$ has density at least $\delta/2$.  Restrict $W$ to the first $k$ letters and
average in the suffix coordinate.  Some fixed suffix $y_0$ leaves density at least $\delta/2$ on a
copy of $[k]^M$.  Proposition~\ref{prop:mdhj} finds an $m$-subspace there.  Replacing its $k$-letter
directions by $(k+1)$-letter directions gives the unique $V$ whose restriction is that subspace.
\end{proof}


\chapter{The density increment dichotomy}
\label{chap:increment}

Fix $k\geq2$, assume $\mathsf{DHJ}_k$, and fix $0<\delta\leq1$.  Following \cite{DKT}, define
\begin{equation}
\label{eq:main-constants}
 m_0=\DHJ(k,\delta/4),\qquad
 \theta=\frac{\delta/4}{(k+1)^{m_0}-k^{m_0}},\qquad
 \eta=\frac{\delta\theta}{48},\qquad
 \gamma=\frac{\delta\eta^2}{k}.
\end{equation}
Every denominator is positive.  We record positivity and the elementary consequences
$\eta<\theta/2$, $\eta\leq\delta/6$, $\gamma\leq\eta^2/2$, and $\gamma\leq3\eta$ as a separate
arithmetic lemma, so later proofs can use them without unfolding the constants.

\begin{lemma}[Numerical facts]
\label{lem:numerical-facts}
\lean{DensityHalesJewett.Parameters.facts}
The constants in \eqref{eq:main-constants} are positive and satisfy all the inequalities listed
above, together with every elementary inequality invoked below.
\end{lemma}

For $m\geq1$ and $\varepsilon>0$, abbreviate
\[
  n(m,\varepsilon)=\left\lceil\varepsilon^{-1}(k+1)^m m\right\rceil.
\]

\section{Correlated fibers}

\begin{lemma}[Every line has correlated fibers]
\label{lem:correlated-fibers}
\lean{DensityHalesJewett.exists_subspace_correlated_fibers}
\uses{lem:uniform-fibers,cor:gr-two-color,lem:threshold-average,lem:lines-equiv,lem:numerical-facts}
Let $m\geq m_0$ and
$n\geq n(\GR(k,m),\eta^2/2)$.  If $A\subseteq[k+1]^n$ has density at least $\delta$, then there
are $l<n$ and an $m$-dimensional subspace $U\subseteq[k+1]^l$ such that
\begin{align*}
 \dens(A_u)&\geq\delta-\eta^2/2 &&\text{for every }u\in U,\\
 \dens\left(\bigcap_{u\in\ell}A_u\right)&\geq\theta
 &&\text{for every }\ell\in\Lines(U\restr[k]).
\end{align*}
\end{lemma}

\begin{proof}
Lemma~\ref{lem:uniform-fibers} first gives a $\GR(k,m)$-subspace $V$ satisfying the pointwise fiber
bound.  Color a line $\ell$ of $V\restr[k]$ good when
$\dens(\bigcap_{v\in\ell}A_v)\geq\theta$.  Corollary~\ref{cor:gr-two-color} gives an
$m$-subspace $Y$ all of whose lines are good or all of whose lines are bad.

The bad alternative is impossible.  Choose an $m_0$-subspace $Z\subseteq Y$.  The average, over
suffixes $y$, of $\dens_Z\{z:z\concat y\in A\}$ is at least $\delta-\eta^2/2\geq\delta/2$.
The threshold estimate gives a set $B$ of suffixes of density at least $\delta/4$ on which this
relative density is at least $\delta/4$.  By the definition of $m_0$, each such slice contains a
line $\ell_y$ of $Z$.  There are exactly $(k+1)^{m_0}-k^{m_0}$ possible lines in $Z$, so one line
$\ell_0$ is selected on a suffix set of density at least
\[
 \frac{\delta/4}{(k+1)^{m_0}-k^{m_0}}=\theta.
\]
Thus $\ell_0$ is good, a contradiction.  Extend $Y$ uniquely to a $(k+1)$-letter subspace $U$.
\end{proof}

\begin{lemma}[Many lines in a dense slice]
\label{lem:many-lines}
\lean{DensityHalesJewett.exists_subspace_many_lines}
\uses{lem:correlated-fibers,lem:threshold-average,lem:numerical-facts}
Under the hypotheses of Lemma~\ref{lem:correlated-fibers}, either some $m$-subspace $X$ satisfies
$\dens_X(A)\geq\delta+\eta^2/2$, or there is an $m$-subspace $W$ with
\[
 \dens_W(A)\geq\delta-2\eta
\]
and
\[
 \card{\{\ell\in\Lines(W\restr[k]):\ell\subseteq A\}}
 \geq\frac\theta2\card{\Lines(W\restr[k])}.
\]
\end{lemma}

\begin{proof}
Assume the first alternative fails and take $U$ from Lemma~\ref{lem:correlated-fibers}.  For a
suffix $y$, let $f(y)=\dens_{U\concat y}(A)$.  Its average is at least
$\delta-\eta^2/2$, while the failure of the first alternative gives
$f(y)<\delta+\eta^2/2$ pointwise.  If the set $H_1=\{y:f(y)\geq\delta-2\eta\}$ had density below
$1-\eta$, splitting the average over $H_1$ and its complement would contradict these two bounds.
Hence $\dens(H_1)\geq1-\eta$.

Let $g(y)$ be the proportion of lines $\ell$ of $U\restr[k]$ for which
$y\in\bigcap_{u\in\ell}A_u$.  Double counting gives $\E_y g(y)\geq\theta$, so the elementary
threshold bound gives $\dens(H_2)\geq\theta/2$ for
$H_2=\{y:g(y)\geq\theta/2\}$.  Since $\eta<\theta/2$, the union bound shows
$H_1\cap H_2\neq\varnothing$.  For $y_0$ in the intersection, take $W=U\concat y_0$.
\end{proof}

\section{Insensitive sets and structured correlation}

\begin{definition}[Insensitive set]
\label{def:insensitive}
\lean{DensityHalesJewett.Insensitive}
Let $i\neq j$ be letters.  Words $x,y\in\alpha^n$ are $(i,j)$-equivalent if they agree at every
coordinate at which the common letter is not $i$ or $j$; equivalently, for each
$a\notin\{i,j\}$ the sets of coordinates carrying $a$ in $x$ and $y$ coincide.  A set $D$ is
$(i,j)$-insensitive if membership is constant on these equivalence classes.
\end{definition}

\begin{lemma}[Boolean operations preserve insensitivity]
\label{lem:insensitive-boolean}
\lean{DensityHalesJewett.Insensitive.compl}
\uses{def:insensitive}
Insensitive subsets are closed under complements, finite unions, and finite intersections.  The
same is true relative to a subspace after transport through its parameterization.
\end{lemma}

Set $\lambda=(k+1)/k$ and choose the integer
\begin{equation}
\label{eq:M0}
 M_0=\max\left\{m_0,
 \left\lceil\frac{\log(\eta^{-1})}{\log\lambda}\right\rceil\right\}.
\end{equation}
The formal definition can avoid real logarithms: take the least $M\geq m_0$ satisfying
$(k/(k+1))^M\leq\eta$.  Existence follows from geometric decay.

\begin{lemma}[A large insensitive intersection]
\label{lem:large-insensitive-intersection}
\lean{DensityHalesJewett.exists_large_insensitive_intersection}
\uses{lem:many-lines,def:insensitive,lem:lines-equiv,lem:numerical-facts}
Let $m\geq M_0$, and suppose $A\subseteq[k+1]^n$ is line-free, has density at least $\delta$, and
has density below $\delta+\eta^2/2$ in every $m$-subspace.  Then some $m$-subspace $W$ contains
$C=\bigcap_{i=1}^k C_i$, with each $C_i$ $(i,k+1)$-insensitive in $W$, such that
\begin{align*}
 \dens_W(C)&\geq\theta/4,\\
 \dens_W(A\cap(W\smallsetminus C))&\geq(\delta+6\eta)\dens_W(W\smallsetminus C),\\
 \dens_W(A\cap(W\smallsetminus C))&\geq\delta-3\eta.
\end{align*}
\end{lemma}

\begin{proof}
Take $W$ from the second alternative of Lemma~\ref{lem:many-lines}.  Every line $\ell$ in
$W\restr[k]$ has a unique extension $\bar\ell$ to the $(k+1)$-letter cube.  Let $B$ consist of
the endpoints $\bar\ell(k+1)$ for which $\ell\subseteq A$, and put
$C=B\cup(A\cap(W\restr[k]))$.

For $x\in W$, let $x^{k+1\to i}$ replace every $k+1$ in its parameter word by $i$, and define
\[
 C_i=\{x:x^{k+1\to i}\in A\cap(W\restr[k])\}.
\]
Then $C_i$ is $(i,k+1)$-insensitive and direct inspection gives $C=\bigcap_iC_i$.  The endpoint
map $\ell\mapsto\bar\ell(k+1)$ is injective, whence
\[
 \card C\geq\frac\theta2\big((k+1)^m-k^m\big)
 \geq\frac\theta4(k+1)^m.
\]
The penultimate inequality follows from $m\geq M_0$ and the definition of $M_0$.

If $x\in A\cap C$ and $x$ uses the letter $k+1$, the definitions supply all the first $k$ points
of the line ending at $x$, while $x$ supplies the last point.  This contradicts line-freeness.
Thus $A\cap C\subseteq W\restr[k]$, and
$\dens_W(A\cap C)\leq(k/(k+1))^m\leq\eta$.  Subtracting this from
$\dens_W(A)\geq\delta-2\eta$ gives the last displayed bound.  Finally divide by
$\dens_W(W\smallsetminus C)\leq1-\theta/4$ and use
$(\delta-3\eta)/(1-\theta/4)\geq\delta+6\eta$, one of the recorded numerical facts.
\end{proof}

\begin{corollary}[Correlation with a structured set]
\label{cor:structured-correlation}
\lean{DensityHalesJewett.exists_structured_correlation}
\uses{lem:large-insensitive-intersection,lem:insensitive-boolean,lem:numerical-facts}
Let $m\geq M_0$ and let $A\subseteq[k+1]^n$ have density at least $\delta$ and contain no line.
Then some $m$-subspace $W$ has subsets $D_1,\ldots,D_k$, with $D_i$ $(i,k+1)$-insensitive in
$W$, such that for $D=\bigcap_iD_i$,
\[
 \dens_W(D)\geq\gamma,
 \qquad
 \dens_W(A\cap D)\geq(\delta+\gamma)\dens_W(D).
\]
\end{corollary}

\begin{proof}
If an $m$-subspace $X$ has density at least $\delta+\eta^2/2$, take $W=X$ and every $D_i=W$.
Otherwise take $C=\bigcap_iC_i$ from Lemma~\ref{lem:large-insensitive-intersection}.  Partition
$W\smallsetminus C$ into
\[
 P_1=W\smallsetminus C_1,
 \qquad
 P_i=(W\smallsetminus C_i)\cap C_1\cap\cdots\cap C_{i-1}\quad(i>1).
\]
Write $\lambda_i=\dens_W(P_i)/\dens_W(W\smallsetminus C)$ and let $\delta_i$ be the relative density
of $A$ in $P_i$ (zero when $P_i$ is empty).  Then
$\sum_i\lambda_i\delta_i\geq\delta+6\eta$.  The total $\lambda$-mass of indices with
$\delta_i\geq\delta+3\eta$ is at least $3\eta$: otherwise the average is strictly below
$\delta+6\eta$, using $\delta_i\leq1$.  Hence some $i_0$ satisfies both
$\lambda_{i_0}\geq3\eta/k$ and $\delta_{i_0}\geq\delta+3\eta$.

Take $D_i=C_i$ for $i<i_0$, $D_{i_0}=W\smallsetminus C_{i_0}$, and $D_i=W$ for $i>i_0$.
Then $D=P_{i_0}$ and Boolean closure gives the required insensitivity.  The second lower bound in
Lemma~\ref{lem:large-insensitive-intersection} yields
\[
 \dens_W(D)\geq(3\eta/k)(\delta-3\eta)\geq\gamma,
\]
while $\delta_{i_0}\geq\delta+3\eta\geq\delta+\gamma$ gives the correlation estimate.
\end{proof}

\section{Tiling structured sets}

Fix $0<\beta\leq1$ and $m\geq1$.  Let
\begin{equation}
\label{eq:tiling-bound}
 M_1=\MDHJ^*(k,m,\beta),\qquad
 F(m,\beta)=\left\lceil
 \beta^{-1}(k+1+m)^{M_1}(k+1)^{M_1-m}M_1
 \right\rceil.
\end{equation}

\begin{lemma}[Tiling one insensitive set]
\label{lem:tile-insensitive}
\lean{DensityHalesJewett.Insensitive.exists_disjoint_subspaces}
\uses{cor:restricted-subspace,lem:fiber-average,lem:threshold-average}
If $n\geq F(m,\beta)$ and $D\subseteq[k+1]^n$ is $(i,k+1)$-insensitive with
$\dens(D)\geq2\beta$, then there is a pairwise disjoint family $\cV$ of $m$-subspaces contained
in $D$ such that
\[
  \dens\left(D\smallsetminus\bigcup\cV\right)<2\beta.
\]
\end{lemma}

\begin{proof}
Put
$\Theta=\beta(k+1+m)^{-M_1}(k+1)^{m-M_1}$.  Fiber $D$ over a terminal block of length $M_1$.
The set of prefixes whose fiber has density at least $\beta$ has density at least $\beta$.
Corollary~\ref{cor:restricted-subspace} supplies, in each such fiber, an $m$-subspace $V_x$ whose
$k$-letter restriction is contained in the fiber.  Insensitivity upgrades this to
$x\concat V_x\subseteq D$.  There are fewer than $(k+1+m)^{M_1}$ generating words for
$m$-subspaces, so one $V_1$ works for a prefix set $S_1$ of density at least
$\beta(k+1+m)^{-M_1}$.  The disjoint family
$\{x\concat V_1:x\in S_1\}$ covers density at least $\Theta$.

If the uncovered part has density below $2\beta$, stop.  Otherwise move to the preceding unused
block of length $M_1$ and repeat.  Although the global remainder need no longer be insensitive,
each section obtained by fixing all later blocks is insensitive: the removed family in a later
block is described by an insensitive prefix set, and complements preserve insensitivity.  Thus the
same argument applies and adds at least $\Theta$ of fresh, pairwise disjoint subspaces.

After $s$ stages the covered density is at least $s\Theta$, so the construction stops within
$\lfloor\Theta^{-1}\rfloor+1$ stages.  The definition of $F$ guarantees at least this many disjoint
coordinate blocks.  Formalization uses a finite induction whose invariant records disjointness,
containment in $D$, the coverage increment, and insensitivity of every remaining section.
\end{proof}

Define iterates
\[
 F^{(1)}(m,\beta)=F(m,\beta),\qquad
 F^{(r+1)}(m,\beta)=F^{(r)}(F(m,\beta),\beta).
\]

\begin{corollary}[Tiling an intersection]
\label{cor:tile-intersection}
\lean{DensityHalesJewett.Insensitive.exists_disjoint_subspaces_iInter}
\uses{lem:tile-insensitive}
Let $1\leq r\leq k$, $n\geq F^{(r)}(m,\beta)$, and let $D_i$ be
$(i,k+1)$-insensitive for $1\leq i\leq r$.  If $D=\bigcap_{i=1}^rD_i$ has density at least
$2r\beta$, then $D$ contains a pairwise disjoint family $\cV$ of $m$-subspaces such that
\[
 \dens\left(D\smallsetminus\bigcup\cV\right)<2r\beta.
\]
\end{corollary}

\begin{proof}
Induct on $r$.  The base case is Lemma~\ref{lem:tile-insensitive}.  For the step, tile
$D'=D_1\cap\cdots\cap D_r$ up to error $2r\beta$ by pairwise disjoint subspaces of dimension
$F(m,\beta)$.  On each large tile $V$, transport $D_{r+1}\cap V$ to its parameter cube.  If its
relative density is at least $2\beta$, apply Lemma~\ref{lem:tile-insensitive} inside $V$; otherwise
discard the tile.

The final error has two disjoint contributions.  Outside the first-stage tiles it is below
$2r\beta$.  Inside each first-stage tile, either the entire contribution of $D_{r+1}$ is below
$2\beta$ of that tile, or the second-stage tiling leaves below $2\beta$ of that tile.  Summing over
the disjoint first-stage tiles contributes below $2\beta$ globally.  The total is below
$2(r+1)\beta$.
\end{proof}

\section{The dichotomy}

\begin{proposition}[Density increment or a line]
\label{prop:density-increment}
\lean{DensityHalesJewett.density_increment}
\uses{cor:structured-correlation,cor:tile-intersection}
Assume $\mathsf{DHJ}_k$.  For every $0<\delta\leq1$ and $d\geq1$ there is $N(k,d,\delta)$ such
that, whenever $n\geq N(k,d,\delta)$ and $A\subseteq[k+1]^n$ has density at least $\delta$, either
$A$ contains a line or some $d$-subspace $V$ satisfies
\[
  \dens_V(A)\geq\delta+\gamma/2,
\]
where $\gamma$ is defined in \eqref{eq:main-constants}.
\end{proposition}

\begin{proof}
Let $\beta=\gamma^2/(4k)$,
$m(d)=\max\{M_0,F^{(k)}(d,\beta)\}$, and
\[
 N(k,d,\delta)=n(\GR(k,m(d)),\eta^2/2).
\]
Assume $A$ is line-free.  Corollary~\ref{cor:structured-correlation}, applied in an
$m(d)$-subspace $W$, gives $D=\bigcap_{i=1}^kD_i$ with
\[
 \dens_W(D)\geq\gamma,
 \qquad
 \dens_W(A\cap D)\geq(\delta+\gamma)\dens_W(D).
\]
Because $2k\beta=\gamma^2/2\leq\gamma$, Corollary~\ref{cor:tile-intersection} tiles all but density
$\gamma^2/2$ of $D$ by pairwise disjoint $d$-subspaces.  If $T$ is their union, then
\begin{align*}
 \dens_W(A\cap T)
 &\geq(\delta+\gamma)\dens_W(D)-\dens_W(D\smallsetminus T)\\
 &\geq(\delta+\gamma/2)\dens_W(T).
\end{align*}
The last step uses $\dens_W(D)\geq\gamma$ and $\dens_W(T)\leq\dens_W(D)$; it will be isolated as a
small real-arithmetic lemma.  Averaging over the disjoint tiles gives one $V$ with
$\dens_V(A)\geq\delta+\gamma/2$.
\end{proof}


\chapter{Completing the density Hales--Jewett theorem}
\label{chap:iteration}

\section{The binary base case}

\begin{lemma}[Binary lines and comparable sets]
\label{lem:binary-line}
\lean{DensityHalesJewett.binary_line_iff_ssubset}
\uses{def:line}
Identify a binary word with the subset of coordinates on which it is $1$.  Two distinct binary
words are the points of a combinatorial line if and only if the corresponding subsets are strictly
comparable by inclusion.
\end{lemma}

\begin{proof}
The variable coordinates are precisely the set difference between the larger and smaller subset;
the fixed $1$ coordinates form their intersection.  Conversely, a strict inclusion has nonempty
difference and therefore determines a line.
\end{proof}

\begin{lemma}[Density Hales--Jewett for two letters]
\label{lem:dhj-two}
\lean{DensityHalesJewett.dhj_two}
\uses{lem:binary-line}
For every $\delta>0$, sufficiently large dense subsets of $[2]^n$ contain a line.
\end{lemma}

\begin{proof}
A line-free family becomes an antichain by Lemma~\ref{lem:binary-line}.  Sperner's theorem bounds
its cardinality by $\binom n{\lfloor n/2\rfloor}$.  Use the standard estimate
\[
 \binom n{\lfloor n/2\rfloor}/2^n\longrightarrow0
\]
(or the elementary bound $\leq C/\sqrt n$ already available in mathlib) and choose $N$ so that
this ratio is below $\delta$ for $n\geq N$.
\end{proof}

\section{Making the increment uniform}

The paper says that Proposition~\ref{prop:density-increment} completes the proof by a standard
iteration.  A formal proof must ensure that the increments do not shrink to zero as the current
density increases.

\begin{lemma}[Monotone parameters]
\label{lem:uniform-gamma}
\lean{DensityHalesJewett.Parameters.gamma_mono_lowerBound}
\uses{lem:numerical-facts}
Fix $\delta_0>0$.  The bounds in \eqref{eq:main-constants} may be chosen so that for every
$\rho\in[\delta_0,1]$ the increment parameter satisfies
\[
  \gamma(k,\rho)\geq\underline\gamma(k,\delta_0)>0.
\]
\end{lemma}

\begin{proof}
Choose $m_0=\DHJ(k,\delta_0/4)$ once and use this possibly nonminimal dimension for every
$\rho\geq\delta_0$.  Then
\[
 \theta(\rho)\geq
 \frac{\delta_0/4}{(k+1)^{m_0}-k^{m_0}},
\]
and the definitions of $\eta$ and $\gamma$ give a positive lower bound obtained by replacing every
occurrence of $\rho$ by $\delta_0$.  Every proof in Chapter~\ref{chap:increment} only uses the
inequalities guaranteed by these lower bounds, so Proposition~\ref{prop:density-increment} remains
valid with increment $\underline\gamma/2$ above the actual density $\rho$.
\end{proof}

\begin{theorem}[Induction on the alphabet]
\label{thm:dhj-induction}
\lean{DensityHalesJewett.density_hales_jewett_fin}
\uses{lem:dhj-two,prop:density-increment,lem:uniform-gamma}
For every $k\geq2$, $\mathsf{DHJ}_k$ holds.
\end{theorem}

\begin{proof}
Induct on $k$, starting with Lemma~\ref{lem:dhj-two}.  Assume $\mathsf{DHJ}_k$ and fix a density
floor $\delta_0>0$.  Let $\varepsilon=\underline\gamma(k,\delta_0)/2$ and choose
$R\in\N$ with $\delta_0+R\varepsilon>1$.

Build required dimensions backwards.  Put $d_R=1$, and for $j=R-1,\ldots,0$ choose $d_j$ large
enough that the uniform form of Proposition~\ref{prop:density-increment}, with target dimension
$d_{j+1}$, applies in every $d_j$-dimensional cube.  Let $n\geq d_0$ and suppose
$A\subseteq[k+1]^n$ is line-free with density at least $\delta_0$.  Repeatedly apply the dichotomy,
transporting the problem through the parameterization of the subspace obtained at the preceding
stage.  After $j$ stages there is a $d_j$-dimensional subspace $V_j$ with
\[
  \dens_{V_j}(A)\geq\delta_0+j\varepsilon.
\]
Line-freeness is inherited by subspaces, so the increment alternative is forced at every stage.
At $j=R$ the displayed density is greater than one, a contradiction.  Hence $A$ contains a line
and $\mathsf{DHJ}_{k+1}$ follows.
\end{proof}

\begin{proof}[Proof of Theorem~\ref{thm:dhj}]
If the alphabet is empty, choose the bound to be one.  In every allowed dimension a diagonal line
exists, and its containment condition is vacuous because it has no evaluations.  If the alphabet
is a singleton, again choose a positive dimension bound; the unique word is in every
positive-density set and the diagonal line works.  For larger alphabets, choose an equivalence with
$\operatorname{Fin}k$, apply
Theorem~\ref{thm:dhj-induction}, and transport the line back using
\texttt{Combinatorics.Line.map} and reindexing.  Finally clear denominators to recover the public
cardinality inequality.
\end{proof}


\chapter{Arithmetic progressions}
\label{chap:ap}

\section{Base-$k$ encoding}

\begin{definition}[Digital encoding]
\label{def:digital-encoding}
\lean{DensityHalesJewett.baseEncode}
For $k,m\geq1$, define
\[
 \operatorname{enc}_{k,m}(x)=\sum_{i=0}^{m-1}x_i k^i,
 \qquad x\in[k]^m,
\]
where letters are represented by $0,\ldots,k-1$.
\end{definition}

\begin{lemma}[Digital equivalence]
\label{lem:digital-equiv}
\lean{DensityHalesJewett.baseEncodeEquiv}
\uses{def:digital-encoding}
The map $\operatorname{enc}_{k,m}$ is a bijection from $[k]^m$ to
$\{0,\ldots,k^m-1\}$.  In particular, translation by $qk^m$ identifies the word cube with the
$q$th consecutive block of length $k^m$ in $\N$.
\end{lemma}

\begin{lemma}[A line encodes an arithmetic progression]
\label{lem:line-to-ap}
\lean{DensityHalesJewett.Line.baseEncode_isArithmeticProgression}
\uses{def:digital-encoding,def:line}
If $\ell$ is a line in $[k]^m$, then
\[
 \operatorname{enc}_{k,m}(\ell(t))=a+t d\qquad(0\leq t<k),
\]
where
\[
 d=\sum_{i\text{ variable in }\ell}k^i>0.
\]
Thus the encoded line is a nonconstant $k$-term arithmetic progression.
\end{lemma}

\begin{proof}
Split the digital sum into fixed and variable coordinates.  Fixed coordinates contribute $a$;
each variable coordinate contributes $t k^i$.  The variable set is nonempty, so $d>0$.
\end{proof}

\begin{lemma}[Dense complete block]
\label{lem:dense-block}
\lean{DensityHalesJewett.exists_dense_digitBlock}
Let $K\geq1$, $\delta>0$, and $N\geq2K/\delta$.  If
$A\subseteq\{0,\ldots,N-1\}$ has $\card A\geq\delta N$, then some complete consecutive block
$\{qK,\ldots,(q+1)K-1\}$ has relative $A$-density at least $\delta/2$.
\end{lemma}

\begin{proof}
The union of the complete $K$-blocks below $N$ omits fewer than $K$ points.  Hence it contains at
least $\delta N-K\geq\delta N/2$ points of $A$.  If every complete block had density below
$\delta/2$, their union would contain fewer than $(\delta/2)N$ points, a contradiction.
\end{proof}

\begin{proof}[Proof of Theorem~\ref{thm:szemeredi}]
The cases $k=0$ and $k=1$ are discharged directly.  For $k\geq2$, let
$m=\DHJ(k,\delta/2)$ and $K=k^m$.  Choose the public bound at least $2K/\delta$.  By
Lemma~\ref{lem:dense-block}, some complete $K$-block has $A$-density at least $\delta/2$.
Pull its intersection with $A$ back through the bijection of Lemma~\ref{lem:digital-equiv} and
apply Theorem~\ref{thm:dhj}.  The resulting line maps, by Lemma~\ref{lem:line-to-ap} and translation
by the block origin, to the required nonconstant arithmetic progression in $A$.
\end{proof}


\chapter{Lean implementation plan}

The declarations should be implemented in layers so that each file has a narrow import surface.

\begin{enumerate}
\item \texttt{Word.lean}: density, coordinate-sum equivalences, fibers, threshold averaging, and
  finite cardinal identities.
\item \texttt{Subspace.lean}: ranges, composition, restriction of alphabets, line equivalences,
  concatenation of subspaces, and transport through reindexing.
\item \texttt{GrahamRothschild.lean}: finite unions, block canonization, and
  Theorem~\ref{thm:graham-rothschild}.
\item \texttt{UniformFibers.lean}: Proposition~\ref{prop:mdhj},
  Lemma~\ref{lem:uniform-fibers}, and Corollary~\ref{cor:restricted-subspace}.
\item \texttt{Insensitive.lean}: the equivalence relation, Boolean closure, relative
  insensitivity, and the tiling lemmas.
\item \texttt{DensityIncrement.lean}: the numerical parameter structure and all results in
  Chapter~\ref{chap:increment}.
\item \texttt{Main.lean}: the binary case, uniform iteration, alphabet transport, and the proof of
  the final density Hales--Jewett challenge declaration.
\item \texttt{ArithmeticProgression.lean}: digital encoding and the final arithmetic-progression
  challenge declaration.
\end{enumerate}

The most useful implementation invariant is to keep subspaces parameterized.  A statement about
the finite set $[V]$ should usually be pulled back to the full parameter cube, proved there, and
pushed forward using injectivity.  This makes relative density, insensitivity, nested subspaces,
and line counts all reduce to reusable facts about finite function types.

\begin{thebibliography}{9}

\bibitem{DKT}
P. Dodos, V. Kanellopoulos, and K. Tyros,
\emph{A simple proof of the density Hales--Jewett theorem},
International Mathematics Research Notices 2014 (2014), no.~12, 3340--3352;
\href{https://arxiv.org/abs/1209.4986}{arXiv:1209.4986}.

\bibitem{GR}
R. L. Graham and B. L. Rothschild,
\emph{Ramsey's theorem for $n$-parameter sets},
Transactions of the American Mathematical Society 159 (1971), 257--292.

\bibitem{HJ}
A. W. Hales and R. I. Jewett,
\emph{Regularity and positional games},
Transactions of the American Mathematical Society 106 (1963), 222--229.

\bibitem{Sperner}
E. Sperner,
\emph{Ein Satz \"uber Untermengen einer endlichen Menge},
Mathematische Zeitschrift 27 (1928), 544--548.

\end{thebibliography}
